% Abstract + §1 Introduction + §2 Related Work
% Source: user-authored LaTeX, integrated as-is so main.tex can \input{}.

\begin{abstract}
LLM agents are increasingly deployed for user-facing sequential decision-making,
where they act on behalf of users whose preferences are captured in structured
profiles. Errors that violate these profiles often go undetected until trajectory
completion. We take a step toward principled error repair in profile-conditioned
agent trajectories. We identify three challenges: (i) existing methods critique
flat trajectories without modeling which latent assumptions caused the failure;
(ii) search- and uncertainty-based approaches localize \emph{where} to intervene
but not \emph{what} to revise, leading to wasteful re-rollouts; (iii) violations
detected only at completion leave forward-only methods unable to recover
otherwise-valid prefixes. To address these issues, we introduce
\textsc{Intro-Specter}, a three-stage framework that formalizes self-correction
as abductive Bayesian inference over an assumption graph. We first construct an
Assumption-DAG over reasoning steps. We then compute a posterior
$p(a_k \mid \mathcal{V})$ over candidate faulty assumptions via LLM-scored
counterfactual repairs. Finally, we apply a Min-Cost Graph-Edit Repair (MCGR)
procedure that identifies the most likely faulty assumption under the posterior
and re-samples only the downstream subgraph while preserving the unaffected
prefix. We evaluate \textsc{Intro-Specter} against seven self-correction baselines
on profile-conditioned agent benchmarks and find it Holm-Bonferroni-significant
on 18 of 45 model$\times$benchmark cells across eight reconstruction benchmarks
plus a real-data fidelity anchor, peaking at $+33.3$pp on
\textsc{Profile-MuSiQue} $\times$ Gemini Flash ($p = 10^{-5}$), while matching
oracle controls at zero token overhead in the controlled synthetic setting.
\end{abstract}

% =====================================================================
% INTRODUCTION
% =====================================================================
\section{Introduction}
\label{sec:intro}

LLM agents such as ReAct~\citep{yao2023react},
Reflexion~\citep{shinn2023reflexion}, and generative agents~\citep{park2023generative}
are increasingly deployed for sequential decision-making on behalf of
users---scheduling meetings, booking flights, ordering food, and managing
personal workflows~\citep{wang2024survey,liu2024agentbench}. In these settings,
the agent operates over a structured user profile that captures preferences
and constraints, and users often trust the agent's decisions without
re-verifying each intermediate step. Because these agents commit to each
action before moving on, errors that emerge midway through a trajectory---such
as selecting a meal that violates a dietary restriction or booking a hotel
that conflicts with a stated preference---are hard to trace back to the
assumptions that caused them, leaving root causes intact and risking
\emph{user-specific hallucinations}: outputs that are internally coherent
but violate the user's own stated constraints.

Recent work confirms that this problem is both real and systematic.
\citet{singh2024travelplanner} show that injecting user models into multi-step
planning degrades hard constraint satisfaction even as personalized plans are
strongly preferred by users. \citet{sun2026personalization} formalize this as
personalization-induced hallucination, showing that user-specific context can
distort factual reasoning through representational entanglement.
\citet{he2025martingale} provide evidence that LLM belief updates during reasoning
can systematically violate Bayesian rationality, raising concerns about
self-correction mechanisms that rely solely on the model judging its own
trajectory. Together, these findings suggest that reliable correction in
user-conditioned settings requires two ingredients that are not jointly
addressed by current methods: backward attribution from observed errors to the
assumptions that caused them, and external grounding against the user profile
rather than LLM self-evaluation alone.

Several approaches have been proposed to improve LLM agent outputs through
iterative refinement and structured search. Self-Refine~\citep{madaan2023selfrefine}
alternates between feedback and refinement steps using the same model.
Reflexion~\citep{shinn2023reflexion} converts environment feedback into verbal
summaries stored as episodic memory. Tree of Thoughts~\citep{yao2023tot}
maintains diverse reasoning alternatives through branching and backtracking,
and Graph of Thoughts~\citep{besta2024got} generalizes to graph transformations
including refinement, aggregation, and merging. Process-level signals such as
UProp~\citep{duan2025uprop} and URSA~\citep{luo2025ursa} attach uncertainty or
reward estimates to individual reasoning steps. While these methods improve
output quality across a range of tasks, few of them maintain an explicit model
of which latent assumptions caused a failure, and existing approaches do not
prescribe how to selectively repair a trajectory once an error is detected.

In this paper, we propose \textsc{Intro-Specter}, a framework that formalizes
agent self-correction as abductive Bayesian inference over an assumption graph.
Rather than searching forward through a space of thoughts or critiquing flat
trajectories, \textsc{Intro-Specter} maintains an explicit directed acyclic
graph of the assumptions underlying each reasoning step---some derived from the
user profile, others inferred by the model. When a user-profile violation is
detected, \textsc{Intro-Specter} treats the error as evidence, computes a
posterior $p(a_k \mid \mathcal{V})$ over ancestor assumptions via LLM-scored
counterfactual repairs, and solves a minimum-expected-cost graph-edit problem
that selects the most likely faulty assumption under the posterior and
re-samples only the downstream subgraph while reusing the unaffected prefix.
This combination of backward attribution, profile-grounded detection, and
selective repair is, to our knowledge, not jointly addressed by prior work.

\paragraph{Contributions.}
\begin{itemize}[leftmargin=*,itemsep=2pt,topsep=2pt]
\item \textbf{Problem formalization.} We define user-specific hallucination as
a profile-grounded trajectory failure caused by unsupported or contradicted
assumptions about user constraints, and formulate self-correction as abductive
posterior inference over an explicit assumption graph.
\item \textbf{Method.} We introduce an Assumption-DAG that records profile-,
tool-, and model-inferred assumptions with provenance, and a posterior
attribution procedure that scores candidate faulty assumptions via LLM-driven
counterfactual repair.
\item \textbf{Selective repair objective.} We formulate repair as a
decision-theoretic objective trading posterior fault probability against
downstream edit cost, with empirically calibrated abstention for
low-confidence cases.
\item \textbf{Evaluation.} We benchmark across personalized QA, planning,
and sequential-agent tasks, against seven self-correction baselines including
ReAct, Tree of Thoughts, and SelfCheckGPT, with attribution labels,
calibration analysis, ablations, and token-cost comparisons.
\end{itemize}

% =====================================================================
% RELATED WORK
% =====================================================================
\section{Related Work}
\label{sec:related}

\paragraph{Reasoning structures and forward self-correction.}
Chain-of-Thought prompting~\citep{wei2022cot} enables LLMs to decompose problems
into intermediate steps. Self-Consistency~\citep{wang2023selfconsistency} samples
multiple chains and selects via majority vote. Tree of Thoughts~\citep{yao2023tot}
maintains a tree of partial reasoning states with LLM-based evaluation and
BFS/DFS search; Graph of Thoughts~\citep{besta2024got} further allows branch
merging for divide-and-conquer strategies. On the correction side,
Self-Refine~\citep{madaan2023selfrefine} iterates generate--critique--refine
loops using the same LLM; Reflexion~\citep{shinn2023reflexion} stores verbal
post-mortems in episodic memory for future attempts; ReAct~\citep{yao2023react}
interleaves reasoning with tool calls and environment feedback. These methods
have substantially advanced LLM reasoning, but correction operates forward: to
our knowledge, none of these methods explicitly attribute an error at step $T$
to a faulty assumption at step $k<T$ via an explicit assumption graph or belief
state.

\paragraph{Hallucination detection.}
SelfCheckGPT~\citep{manakul2023selfcheckgpt} samples $N$ alternative completions
and checks per-claim consistency for hallucination detection in a zero-resource
setting. HaloScope~\citep{du2024haloscope} learns hallucination scores from
internal model representations. These methods establish that detection is
tractable, but our \textsc{Detection-only} ablation (Section~\ref{sec:results})
shows detection alone produces zero downstream success lift without paired
repair, motivating our three-layer architecture.

\paragraph{Process supervision and uncertainty.}
UProp~\citep{duan2025uprop} decomposes per-step uncertainty into intrinsic
and extrinsic components, showing that uncertainty inherited from earlier
decisions dominates in later steps. URSA~\citep{luo2025ursa} trains a process
reward model for multimodal reasoning. Both establish that error propagation
is a first-class concern, but UProp is diagnostic only and URSA operates at
training time; neither performs trajectory-level repair in user-conditioned
settings.

\paragraph{Personalized agents and user-specific hallucination.}
TravelPlanner+~\citep{singh2024travelplanner} benchmarks personal LLM agents
and shows that personalization can degrade hard constraint satisfaction.
\citet{sun2026personalization} explain a related effect via representational
entanglement between personalization and factual knowledge.
\citet{he2025martingale} provide evidence that LLM belief updates can
systematically entrench priors rather than revising them. These results show
that personalization can corrupt reasoning and that external grounding helps,
but they do not couple user-profile grounding with backward attribution and
trajectory-level repair.

\paragraph{Gap statement.}
Taken together, existing work provides forward reflection, structured search,
hallucination detection, process-level uncertainty, and personalization
analysis. We are not aware of prior work that jointly (i) represents
user-profile assumptions as explicit trajectory nodes with provenance,
(ii) infers a posterior over candidate faulty assumptions after a
profile-grounded violation is detected, and (iii) selectively repairs only
the downstream subgraph while preserving valid prefixes.
\textsc{Intro-Specter} targets this joint problem of detection, attribution,
and selective repair in user-conditioned agent trajectories.
